\hypersetup{
    pdftitle=Digitale Signalverarbeitung,
    pdfauthor=Manfred Storhmann,
}

\title{Digitale Signalverarbeitung}

\author{
    Manfred Strohrmann%
    \thanks{
        \href{mailto:manfred.strohrmann@h-ka.de}{manfred.strohrmann@h-ka.de}
    }
}
\date{\today}

\ifpdf
    \titlehead{
        \includesvg[height=3cm]{HKA_EIT_Logo_Gesamt-h_RGB.svg}
    }
    \subject{Signale und Systeme: Part 2}

    \subtitle{Skript zur Vorlesung}

    \publishers{
        Herausgegeben von:\\
        Manfred Strohrmann
    }

    \newcommand{\SetupTitleBack}{
        \uppertitleback{
            14. Ausgabe, Version 14.0.1, Stand: 2024-05-01
            \medskip

            Dieses Dokument und begleitende Materialien finden Sie unter: \\
            \url{https://www.h-ka.de/eit/Digitale-Signalverarbeitung}
        }
        \lowertitleback{
            Copyright \textcopyright\ 2010 Manfred Strohrmann,
            \medskip

            Published by Manfred Strohrmann. Alle Rechte vorbehalten.
            \medskip

            Dieses Dokument wird als Lehrmaterial für die Vorlesung Digitale Signalverarbeitung an der Hochschule Karlsruhe verwendet. Das Dokument wird unter der folgenden Lizenz veröffentlicht: CC BY-NC-ND 4.0 (\url{https://creativecommons.org/licenses/by-nc-nd/4.0/}).
        }
        % \dedication{
        %     \itshape
        %     Dem akademischen Nachwuchs gewidmet, der die Filter von morgen entwirft.
        % }
    }

    % titleback options are only allowed in twoside mode,
    % so we detect twoside or semitwoside (from KOMA) and enable titleback accordingly.
    \makeatletter
    \newif\ifshowtitleback
    \if@twoside
        \showtitlebacktrue        
    \fi
    \@ifundefined{if@semitwoside}{}{%
        \if@semitwoside
            \showtitlebacktrue
        \fi
    }
    \ifshowtitleback
        \SetupTitleBack
    \fi
    \makeatother

    
\fi